\section{Methodology}
\label{sec:method}

In this section, we present \textbf{Sparse Low-Rank Dynamic Merging (SLD-Merge)}. The core design goal of SLD-Merge is to enable input-adaptive dynamic model merging while completely eliminating batch-dependency and heterogeneity collapse. As illustrated in Figure~\ref{fig:method_overview}, SLD-Merge operates in three phases: (1) offline SVD task-vector decomposition, (2) bounded cosine-similarity routing, and (3) sample-wise parallel execution through Top-1 sparse low-rank adapters.

\begin{figure*}[t]
\centering
\includegraphics[width=0.95\textwidth]{../results/task_wise_performance_b64.png}
\caption{\textbf{Overview of SLD-Merge Pipeline.} \textbf{Left:} In the offline phase, we compute dense task vectors from specialized experts and perform Singular Value Decomposition (SVD) to project them into compact low-rank matrices ($B_k$ and $A_k$). \textbf{Right:} At inference time, input activations are spatially pooled to form $z(x)$. We compute the cosine-similarity against task routing bases $\Phi_k$, apply Top-1 hard selection to produce sparse coefficients $\alpha_k$, and execute a fully parallelized, batch-independent low-rank forward pass.}
\label{fig:method_overview}
\end{figure*}

\subsection{Offline Task Vector SVD Factorization}
Let $W_{base}^{(l)} \in \mathbb{R}^{D_{out} \times D_{in}}$ be the pre-trained base model weight matrix at layer $l$, and let $W_k^{(l)}$ be the corresponding weight matrix of an expert model fully fine-tuned on task $k \in \{1, \dots, K\}$. The dense task vector representing the weight-space parameter shift for task $k$ is defined as:
\begin{equation}
    V_k^{(l)} = W_k^{(l)} - W_{base}^{(l)}
\end{equation}
Instead of storing and loading full-rank task vectors for all $K$ experts—which leads to prohibitive $O(K \cdot N)$ storage scaling—we perform Singular Value Decomposition (SVD) on each $V_k^{(l)}$ offline:
\begin{equation}
    V_k^{(l)} = U_k^{(l)} \Sigma_k^{(l)} ({V'_k}^{(l)})^T
\end{equation}
where $U_k^{(l)} \in \mathbb{R}^{D_{out} \times D_{out}}$ and ${V'_k}^{(l)} \in \mathbb{R}^{D_{in} \times D_{in}}$ are orthogonal matrices, and $\Sigma_k^{(l)}$ contains the singular values. To compress the parameter footprint, we truncate this decomposition to a pre-defined target rank $r \ll \min(D_{out}, D_{in})$ (e.g., $r \in \{4, 8, 16\}$). We construct the lightweight low-rank matrices $B_k^{(l)} \in \mathbb{R}^{D_{out} \times r}$ and $A_k^{(l)} \in \mathbb{R}^{r \times D_{in}}$ by splitting the singular values equally:
\begin{align}
    B_k^{(l)} &= U_k^{(l)}[:, :r] \sqrt{\Sigma_k^{(l)}[:r]} \\
    A_k^{(l)} &= \sqrt{\Sigma_k^{(l)}[:r]} \left( {V'_k}^{(l)}[:, :r] \right)^T
\end{align}
This formulation ensures that $V_k^{(l)} \approx B_k^{(l)} A_k^{(l)}$ with minimal reconstruction error under the Frobenius norm. This decomposition is performed offline once, reducing the additional task-specific storage overhead for $K$ experts by over \textbf{92.5\%}.

\subsection{Bounded Cosine-Similarity Router}
Let $X \in \mathbb{R}^{B \times N \times D}$ be the input activation tensor to layer $l$, where $B$ is the batch size, $N$ is the sequence (or token) length, and $D$ is the embedding dimension. To obtain a global, sequence-level sample representation $z(x)_b \in \mathbb{R}^D$ for each individual sample $b \in \{1, \dots, B\}$, we compute the spatial average over tokens:
\begin{equation}
    z(x)_b = \frac{1}{N} \sum_{n=1}^N X_{b, n, :}
\end{equation}
We define a set of task routing basis vectors $\Phi_k^{(l)} \in \mathbb{R}^D$ representing the signature feature space of each expert $k$. For each sample $b$ and expert $k$, we compute the cosine-similarity routing score:
\begin{equation}
    s_{k, b}^{(l)} = \frac{\langle z(x)_b, \Phi_k^{(l)} \rangle}{\|z(x)_b\|_2 \|\Phi_k^{(l)}\|_2 + \epsilon}
\end{equation}
where $\epsilon = 10^{-8}$ is a numerical stabilizer. 

Crucially, mapping the activation representation onto a bounded, low-dimensional spherical cosine space rather than using unconstrained linear projections provides two key pragmatic advantages: (1) it strictly bounds the routing scores in $[-1, 1]$, suppressing high-frequency activation noise and preventing representation collapse, and (2) it serves as a powerful regularizer, enabling zero-shot and extremely robust few-shot calibration.

\subsection{Top-1 Sparse Gating and Parallel Forward Pass}
Rather than using a soft, dense combination of experts (which would require executing $K$ low-rank forward passes per sample), we enforce extreme computational efficiency by performing Top-1 sparse routing (hard gating). For each sample $b$, we select only the single best-aligned expert task vector:
\begin{equation}
    \hat{k}_b^{(l)} = \arg\max_{k \in \{1, \dots, K\}} s_{k, b}^{(l)}
\end{equation}
The sample-wise routing coefficient is expressed as a sparse one-hot vector:
\begin{equation}
    \alpha_{k, b}^{(l)} = \begin{cases} 1 & \text{if } k = \hat{k}_b^{(l)} \\ 0 & \text{otherwise} \end{cases}
\end{equation}
The output activation $Y_b \in \mathbb{R}^{N \times D}$ of layer $l$ for sample $b$ is then computed as:
\begin{equation}
    Y_b = X_b W_{base}^{(l)} + \alpha_{\hat{k}_b, b}^{(l)} \cdot \left( (X_b A_{\hat{k}_b}^{(l)}) B_{\hat{k}_b}^{(l)} \right)
\end{equation}
where $X_b \in \mathbb{R}^{N \times D}$. To execute this efficiently on modern hardware (e.g., GPUs and edge NPUs), we implement a fully vectorized, parallel forward pass over the batch:
\begin{equation}
    Y = X W_{base}^{(l)} + \sum_{k=1}^K \alpha_k \odot \left( (X A_k^{(l)}) B_k^{(l)} \right)
\end{equation}
where $\alpha_k \in \mathbb{R}^{B \times 1 \times 1}$ is the broadcasted routing coefficient tensor for expert $k$. 

This implementation is mathematically equivalent to processing each sample in complete isolation. Because the routing coefficient $\alpha_{k,b}$ is computed sample-by-sample and applied element-wise, there is zero averaging across the batch dimension. Consequently, SLD-Merge is completely batch-independent: the output representation and downstream prediction for any sample $b$ are identical regardless of whether the sample is evaluated alone ($B=1$) or within a massive, highly heterogeneous batch.

\subsection{Activation-Space Mean Initialization}
A significant challenge with hard Top-1 routing is that the $\arg\max$ operation is non-differentiable, preventing direct gradient-based optimization of the routing basis vectors $\Phi_k^{(l)}$ via backpropagation. While prior works employ complex and unstable reinforcement learning, Gumbel-Softmax relaxations, or straight-through estimators, we propose an elegant and highly pragmatic alternative: \textbf{Activation-Space Mean Initialization}.

We set each routing basis vector $\Phi_k^{(l)}$ to the empirical mean activation of its corresponding task $k$, computed over a small, unlabeled calibration split $D_{cal, k}$ (e.g., 128 samples per task):
\begin{equation}
    \Phi_k^{(l)} = \frac{1}{|D_{cal, k}|} \sum_{x \in D_{cal, k}} z(x)^{(l)}
\end{equation}
This initialization aligns each task basis vector exactly with the representative feature centroid of that task's activations at layer $l$. 
As we demonstrate in our experimental evaluation, this simple, zero-shot initialization achieves peak joint multi-task accuracy completely training-free—requiring zero backpropagation steps during calibration, which makes it incredibly simple to deploy in streaming production environments.

\subsection{Autonomous Classification Head Selection}
In multi-task networks created from individual expert models, each task expert retains its own task-specific linear classification head (e.g., mapping shared representations to 10 classes). While prior dynamic weight merging works employ an oracle at inference time to select the classification head using ground-truth labels, this introduces information leakage and limits deployment autonomy. 

To achieve true, stateless autonomy, SLD-Merge introduces a sample-wise, layer-averaged classification head selection rule. For each incoming sample $b$, we compute the average routing score across all specialized blocks $\mathcal{L}' = \{9, 10, 11\}$:
\begin{equation}
    \bar{s}_{k,b} = \frac{1}{|\mathcal{L}'|} \sum_{l \in \mathcal{L}'} s_{k,b}^{(l)}
\end{equation}
We then dynamically route the sample's pooled intermediate representation to classification head $\hat{k}_b$, defined completely autonomously as:
\begin{equation}
    \hat{k}_b = \arg\max_{k \in \{1, \dots, K\}} \bar{s}_{k,b}
\end{equation}
Because the visual representation spaces of our four tasks (MNIST, FashionMNIST, CIFAR-10, SVHN) are highly separable, our cosine router distinguishes incoming domains with $\approx 100\%$ accuracy, ensuring that our fully autonomous pipeline achieves identical performance to oracle head selection without any ground-truth leakage.
